% Figure 7 draft A: triangle sites; a grid of pictures: each mode, its t image, its b image; brackets name A1 and E.
\documentclass[10pt,border=1mm]{standalone}
\usepackage[T1]{fontenc}
\usepackage{figurestyle}
\input{primitives.tex}
\input{numbers.tex}
% shared drawing for the figure 7 drafts: a mode picture is three discs on the three version sites;
% disc area = |amplitude|, ink positive, white negative, a dash for zero. Sites at 90, 210, 330 degrees.
% \modepic{x}{y}{R}{a}{b}{c}{style}  style: circle (torsion circle) or tri (triangle)
\newcommand\modepic[7]{\begin{scope}[shift={(#1,#2)}]
  \ifnum#7=1 \draw[fine] (0,0) circle[radius=#3]; \else \draw[fine] (90:#3)--(210:#3)--(330:#3)--cycle; \fi
  \foreach \mang/\mc in {90/#4,210/#5,330/#6} {
    \pgfmathsetmacro{\r}{.42*#3*sqrt(abs(\mc))}\pgfmathsetmacro{\pos}{\mc>0?1:0}
    \ifdim\r pt>0.6pt \ifnum\pos=1 \fill[PlateInk] (\mang:#3) circle[radius=\r]; \else \draw[fine,fill=white] (\mang:#3) circle[radius=\r]; \fi
    \else \draw[fine] ($(\mang:#3)+(-.07,0)$)--($(\mang:#3)+(.07,0)$); \fi}
  \end{scope}}
% amplitudes: v1=(1,1,1)/sqrt3; v2=(2,-1,-1)/sqrt6; v3=(0,1,-1)/sqrt2; t(a,b,c)=(c,a,b); b(a,b,c)=(a,c,b)
\def\sA{.5774}\def\sB{.8165}\def\sC{.4082}\def\sD{.7071}

\begin{document}
\begin{tikzpicture}[plate]
\path[use as bounding box] (0,0) rectangle (15,8.4);
% key: the three sites, t cycles, b fixes H and swaps the others
\begin{scope}[shift={(2.6,4.6)}]
  \pgfmathsetmacro{\R}{1.5}
  \draw[fine] (90:\R)--(210:\R)--(330:\R)--cycle;
  \draw[tpath] (100:{\R+.25}) arc[start angle=100,end angle=200,radius={\R+.25}];
  \draw[tpath] (220:{\R+.25}) arc[start angle=220,end angle=320,radius={\R+.25}];
  \draw[tpath] (340:{\R+.25}) arc[start angle=340,end angle=440,radius={\R+.25}];
  \draw[bpath,<->] (210:{\R-.35})--(330:{\R-.35});
  \draw[blift] (90:{\R+.25}) arc[start angle=-40,end angle=220,radius=.3];
  \foreach \a/\l in {90/H,210/tH,330/t^2H} \node[circle,draw=PlateInk,fill=white,line width=.55pt,inner sep=1pt,minimum size=6.5mm] at (\a:\R) {$\l$};
  \node[lab] at (150:{\R+.55}) {$t$}; \node[lab] at (270:{\R+.55}) {$t$}; \node[lab] at (30:{\R+.55}) {$t$};
  \node[lab] at (270:{\R-.75}) {$b$}; \node[lab,anchor=south] at (90:{\R+.72}) {$b$};
\end{scope}
% the grid
\foreach \col/\x/\a/\b/\c/\ta/\tb/\tc/\ba/\bb/\bc in {1/7.0/\sA/\sA/\sA/\sA/\sA/\sA/\sA/\sA/\sA, 2/10.0/\sB/-\sC/-\sC/-\sC/\sB/-\sC/\sB/-\sC/-\sC, 3/13.0/0/\sD/-\sD/-\sD/0/\sD/0/-\sD/\sD} {
  \modepic{\x}{6.6}{.62}{\a}{\b}{\c}{0}
  \modepic{\x}{4.4}{.62}{\ta}{\tb}{\tc}{0}
  \modepic{\x}{2.2}{.62}{\ba}{\bb}{\bc}{0}
  \node[lab] at (\x,7.55) {$v_\col$};}
\node[lab,anchor=east] at (5.9,6.6) {$v$}; \node[lab,anchor=east] at (5.9,4.4) {$t\,v$}; \node[lab,anchor=east] at (5.9,2.2) {$b\,v$};
\draw[fine] (6.2,8.0)--(6.2,7.85)--(7.8,7.85)--(7.8,8.0); \node[lab,anchor=south] at (7.0,8.02) {$\mathrm A_1$};
\draw[fine] (9.2,8.0)--(9.2,7.85)--(13.8,7.85)--(13.8,8.0); \node[lab,anchor=south] at (11.5,8.02) {$\mathrm E$};
\node[sub,anchor=north] at (10.0,1.3) {disc area is the amplitude; ink positive, white negative, a dash for zero};
\end{tikzpicture}
\end{document}
